%==============================
\section{Simultaneous Planning for Communication and Collaboration}
\label{ch05:sec:ch05_interm_team}

In this section, a multi-agent system is considered in which complex tasks with
temporal constraints are released online. The focus is on an online co-optimization problem in
which collaborative task execution and communication decisions are synthesized
simultaneously. In contrast to periodic team-wide communication or intermittent
pairwise communication schemes, the proposed framework determines
\emph{when}, \emph{where}, and \emph{how} communication or collaboration
should occur according to the task information available in real time. To this
end, a branch-and-bound framework is developed in which task planning and
communication events are encoded jointly within a unified search structure. The
resulting objective balances task-execution efficiency against communication
latency, thereby improving robustness under spatially and temporally varying
task distributions. In addition, communication events are refined through an
iterative procedure that selects suitable communication locations and timings.
The effectiveness of the overall method is demonstrated in large-scale
simulations and hardware experiments against several strong baselines.

The technical development of this section proceeds as follows.
Section~\ref{ch05:sec:problem} formulates the workspace model, the online task
model, the communication constraints, and the associated optimization problem.
Section~\ref{ch05:subsec:coord-framework} then introduces an adaptive
coordination framework in which a new objective function supports robust
online decision making under uncertain task arrivals and dynamic operating
conditions. Next, Section~\ref{ch05:subsec:comm-event} addresses the
optimization of communication events by determining appropriate communication
locations and timings. Finally, Section~\ref{ch05:subsec:online} presents the
overall execution logic and analyzes the computational properties of the
proposed method.







%==============================
%==============================
\subsection{Collaborative fleet under communication constraints}
\label{ch05:sec:problem}

This subsection formulates the fleet model, the online collaborative-task
model, and the communication constraints used throughout
Section~\ref{ch05:sec:ch05_interm_team}. In analogy with the robot-team
formulations introduced in Chapter~\ref{chap:bottom-up} and
Chapter~\ref{chap:top-down}, the starting point is again a collection of
robots evolving in a common workspace and executing timed plans. The present
setting differs in two important respects. First, collaborative tasks are
released online and must be completed under temporal relations that are only
partially known in advance. Second, global information exchange cannot be
maintained continuously and must itself be coordinated through communication
events. Consequently, communication is not an auxiliary mechanism added after
task planning, but an integral part of the planning model.

\subsubsection{Workspace and robot model}
\label{ch05:subsec:env-task-model}

Following the local-execution viewpoint in
Chapter~\ref{chap:bottom-up} and the team-level viewpoint in
Chapter~\ref{chap:top-down}, consider a fleet of $N$ robots operating in a
common workspace $\mathcal{W} \subset \mathbb{R}^2$. Let
$\mathcal{N} \triangleq \{1, \cdots, N\}$ denote the robot index set. For each
robot $i \in \mathcal{N}$, let $\mathcal{W}_i \subseteq \mathcal{W}$ be the
admissible workspace, let $p_i(t) \in \mathcal{W}_i$ be its position at time
$t \geq 0$, let $d_i > 0$ be its sensing radius, and let
$\mathcal{A}_i$ be its set of primitive actions. The action set
$\mathcal{A}_i$ reflects the heterogeneous capabilities of robot $i$. For
example, repair and rescue actions require dexterous manipulation, whereas
delivery actions require transportation capability. The motion of robot $i$ is
described by its velocity $v_i(t) \in \mathbb{R}^2$, which satisfies the bound
$\|v_i(t)\| \leq v_i^{\texttt{max}}$ for all $t \geq 0$.
For each robot $i \in \mathcal{N}$, a nominal local plan is a timed sequence
$\tau_i \triangleq (t_i^1, g_i^1, a_i^1)(t_i^2, g_i^2, a_i^2)\cdots$,
where, for every $\ell \geq 1$, the scalar $t_i^\ell \geq 0$ is the execution
time, the point $g_i^\ell \in \mathcal{W}_i$ is the goal position, and the
symbol $a_i^\ell \in \mathcal{A}_i$ is the action to be performed. Under
$\tau_i$, robot $i$ moves toward $g_i^\ell$ and starts action $a_i^\ell$ at
time $t_i^\ell$. The nominal collective plan is
$J \triangleq \{\tau_i\}_{i \in \mathcal{N}}$.
This plan representation plays the same conceptual role as the execution-level
objects used in Chapter~\ref{chap:bottom-up} and the team-level plan objects
used in Chapter~\ref{chap:top-down}. The difference is that the present plan
must also accommodate future communication decisions and online collaborative
tasks.

\subsubsection{Online collaborative temporal tasks}
\label{ch05:subsec:task-spec}

In analogy with the task formulations developed in the previous chapters, the
present model still captures collaborative activities together with temporal
relations among them. The main difference is that the task set is not fixed in
advance. In Chapter~\ref{chap:bottom-up}, the tasks are assigned locally and
planned from the outset, whereas in Chapter~\ref{chap:top-down}, the mission is
specified globally before decomposition and assignment. Here, by contrast,
collaborative tasks are released online and become known only after detection
by the fleet.
Collaborative tasks are released online with unknown, and possibly time-varying,
spatio-temporal distributions. Each task is represented by
$
\omega_m \triangleq
\bigl(S_m, \eta_m, \{(n_j, a_j)\}_{j=1}^{J_m}\bigr),
$
where $S_m \subseteq \mathcal{W}$ is the task region, $\eta_m > 0$ is the
required execution duration, and each pair $(n_j, a_j)$ specifies that $n_j$
robots are required to perform action $a_j$ within the region $S_m$.
A task becomes known only after it is detected by the fleet. Detection occurs
when at least one robot can sense the task region. More precisely, task
$\omega_m$ is detected at time $t_m$ if there exists a robot
$i \in \mathcal{N}$ such that the distance from $p_i(t_m)$ to $S_m$ is within
the sensing range $d_i$, and the line of sight from $p_i(t_m)$ to $S_m$ is not
blocked by an obstacle. Tasks detected locally are later shared through the
team-wise communication process introduced in the sequel.

\begin{definition}
The set of tasks detected by time $t \geq 0$ is
$\Omega_t \triangleq \{\omega_m \mid t_m \leq t\}$.
For each task $\omega_m \in \Omega_t$, let $t_m^{\texttt{s}} > 0$ and
$t_m^{\texttt{f}} > 0$ denote its start and finish times, respectively.
\end{definition}

As in the task-coordination models developed in the previous chapters, explicit
temporal relations are imposed between task intervals. For any two tasks
$\omega_p, \omega_q \in \Omega_t$, the precedence relation
$\omega_p \preceq \omega_q$ requires
$t_p^{\texttt{f}} \leq t_q^{\texttt{s}}$. The mutual-exclusion relation
$\omega_p \parallel \omega_q$ requires the intervals
$[t_p^{\texttt{s}}, t_p^{\texttt{f}}]$ and
$[t_q^{\texttt{s}}, t_q^{\texttt{f}}]$ to be disjoint, without imposing an
order between them. The concurrency relation $\omega_p \sim \omega_q$ permits
the two intervals to overlap. The set of all temporal relations known by time
$t$ is given by:
\begin{equation}
\mathcal{P}_t \triangleq
\left\{
(\omega_i, \omega_j, \kappa)
\;\middle|\;
\omega_i, \omega_j \in \Omega_t,\;
\kappa \in \{\preceq, \parallel, \sim\}
\right\},
\label{ch05:eq:task-constraints}
\end{equation}
where the relation set in \eqref{ch05:eq:task-constraints} becomes available
only after the corresponding tasks have been detected.
A task $\omega_m \in \Omega_t$ is \emph{completed} at time $t_m^{\texttt{f}}$ if the
following two conditions hold. First, at time $t_m^{\texttt{s}}$, enough robots
arrive at $S_m$ and, for each required pair $(n_j, a_j)$, a team of $n_j$
robots continuously performs action $a_j$ within $S_m$ for duration $\eta_m$.
Second, the interval $[t_m^{\texttt{s}}, t_m^{\texttt{f}}]$ satisfies all
relations in $\mathcal{P}_t$ that involve $\omega_m$.
For more complex temporal specifications, the relation set
$\mathcal{P}_t$ may be derived from the R-poset construction in
Section~\ref{ch04:subsec:task-decompose-posets}. In particular,
Definition~\ref{ch04:def:poset} and \eqref{ch04:eq:poset} encode precedence and
mutual-exclusion relations in a more general partial-order form, of which the
relations in \eqref{ch05:eq:task-constraints} are a task-level specialization.



\subsubsection{Communication constraints}
\label{ch05:subsec:comm-const}

The fleet cannot maintain continuous real-time communication because of physical
limitations in signal propagation and environmental occlusion. Successful communication between
robots $i, j \in \mathcal{N}$ requires the communication quality $Q_{ij}(t)$ to
exceed a threshold $\delta > 0$. The communication quality is modeled by:
\begin{equation}
Q_{ij}(t) \triangleq
P_t - PL\bigl(d_{ij}(t)\bigr) - \alpha d_{ij}^{\mathrm{obs}}(t),
\label{ch05:eq:comm-quality}
\end{equation}
where $P_t$ is the transmission power in dB,
$d_{ij}(t) \triangleq \|p_i(t) - p_j(t)\|$ is the distance between robots $i$
and $j$, $PL(d_{ij}(t))$ is the free-space path-loss term, $d_0$ is a reference
distance, $n$ is the path-loss exponent, $d_{ij}^{\mathrm{obs}}(t)$ is the
total obstacle length along the straight segment joining $p_i(t)$ and $p_j(t)$,
and $\alpha > 0$ is the attenuation coefficient. More precisely, the path-loss
term is given by:
\begin{equation}
PL\bigl(d_{ij}(t)\bigr) \triangleq
PL(d_0) + 10 n \log_{10}\!\left(\frac{d_{ij}(t)}{d_0}\right),
\label{ch05:eq:path-loss}
\end{equation}
according to which these communication constraints induce a time-varying graph
$G(t) \triangleq (\mathcal{N}, \mathcal{E}(t))$, where
$\mathcal{E}(t) \triangleq
\{e_{ij} \mid Q_{ij}(t) > \delta\}$.
Let $\mathcal{G}_{\texttt{c}}$ denote the set of connected graphs over
$\mathcal{N}$. The fleet achieves team-wise global communication at time
$t > 0$ if $G(t) \in \mathcal{G}_{\texttt{c}}$. This communication model is
intermittent at the team level. Global connectivity is required only at selected
communication events, so that task information and local plans can be merged and
updated in a centralized-style manner when such events occur.


A communication event is the collection
$\mathbf{C} \triangleq \{\mathbf{c}_i\}_{i \in \mathcal{N}}$,
where $\mathbf{c}_i \triangleq (t_i^{\texttt{c}}, p_i^{\texttt{c}})$ denotes
the communication time and communication position assigned to robot $i$. The
team-wise event occurs when all robots reach their designated communication
positions at a common time $t^{\texttt{c}}$, that is,
$t^{\texttt{c}} = t_i^{\texttt{c}}$ for all $i \in \mathcal{N}$, and the graph
satisfies
\begin{equation}
G(t^{\texttt{c}}) \in \mathcal{G}_{\texttt{c}},
\label{ch05:eq:comm-graph}
\end{equation}
where $\mathcal{G}_{\texttt{c}}$ is the set of connected communication graphs
over the robot set $\mathcal{N}$.
When \eqref{ch05:eq:comm-graph} holds, the robots exchange critical information,
including locally detected tasks, local plans, and future communication events.
Accordingly, each nominal local plan $\tau_i$ should be understood as containing both
task-execution decisions and communication decisions.
For a nominal collective plan $J$ and a horizon $T > 0$, let
$N_T(J) \triangleq
\left|\left\{\omega_m \in \Omega_T\;\middle|\;t_m^{\texttt{f}}(J) \leq T\right\}
\right|$
be the number of tasks completed by time $T$ under $J$. The associated long-run
average completion rate is
\begin{equation}\label{ch05:eq:objective}
  \eta(J) \triangleq \lim_{T \to \infty} \frac{N_T(J)}{T}
\end{equation}
where the objective is to maximize $\eta(J)$ subject to the temporal constraints in
\eqref{ch05:eq:task-constraints} and the communication-connectivity condition
in \eqref{ch05:eq:comm-graph}.


\begin{problem}
\label{ch05:prob:collective-plan}
Given the robot fleet in
Section~\ref{ch05:subsec:env-task-model}, the online collaborative-task model
in Section~\ref{ch05:subsec:task-spec}, and the communication constraints in
Section~\ref{ch05:subsec:comm-const}, synthesize the collective plan
$J \triangleq \{\tau_i\}_{i \in \mathcal{N}}$ so as to maximize the long-run
average task-completion rate in \eqref{ch05:eq:objective}, subject to the
temporal relation constraints in \eqref{ch05:eq:task-constraints} and the
team-wise communication condition in \eqref{ch05:eq:comm-graph}.
\end{problem}





















% ==============================================================
% ==============================================================
\subsection{Adaptive coordination for communication and collaboration}
\label{ch05:subsec:coord-framework}

This subsection develops the replanning layer that couples task selection with scheduled communication restoration. It first introduces the horizon-wise objective used to score candidate collective plans and then presents the branch-and-bound procedure used to optimize task allocation and the next communication event jointly.

\subsubsection{Design of adaptive optimization objective}
\label{ch05:subsubsec:new-obj}

To address unknown spatio-temporal task distributions while still pursuing the
global objective in \eqref{ch05:eq:objective}, an adaptive optimization
objective is introduced for each execution horizon. The coordination framework
is triggered at every communication event at time $t$. Based on the currently
detected task set $\Omega_t$, the framework jointly optimizes the future local
plans $J \triangleq \{\tau_i\}_{i \in \mathcal{N}}$ together with the next
communication event. Since the team-wise intermittent communication protocol
restores global information exchange at each event, the tasks assigned in the
interval $[t, t^{\texttt{c}}]$ are completed before the next replanning phase
at time $t^{\texttt{c}}$. The long-run rate $\eta(J)$ in
\eqref{ch05:eq:objective} is therefore approximated locally by the task
completion rate over the interval $[t, t^{\texttt{c}}]$.

At a communication event at time $t$, let $\omega_i^\ell$ denote the $\ell$th
task assigned to robot $i$ after time $t$ under the candidate plan $\tau_i$.
Let $t_{i,\ell}^{\texttt{s}}$ and $t_{i,\ell}^{\texttt{f}}$ denote the
corresponding start and finish times. Let $\mathcal{C}_t$ denote the
feasibility constraints induced by the temporal relations in
\eqref{ch05:eq:task-constraints}.
The adaptive optimization problem is formulated as:
\begin{equation}
\begin{aligned}
\max_{\{\tau_i\},\, t^{\texttt{c}}} \quad &
\frac{
\left|
\left\{
\omega_m \in \Omega_t \mid t_m^{\texttt{f}} \leq t^{\texttt{c}}
\right\}
\right|
}{
t^{\texttt{c}} - t
},
\\
\text{s.t.} \quad &
G(t^{\texttt{c}}) \in \mathcal{G}_{\texttt{c}},
\quad \mathcal{C}_t \text{ holds},
\\
&
t < t_{i,\ell}^{\texttt{f}} \leq t^{\texttt{c}},
\quad \forall i \in \mathcal{N},
\\
&
t_{i,\ell+1}^{\texttt{s}}
\geq t_{i,\ell}^{\texttt{f}}
+ T_{\text{travel}}(\omega_i^\ell, \omega_i^{\ell+1}),
\quad \forall i \in \mathcal{N},
\end{aligned}
\label{ch05:eq:adaptive_objective}
\end{equation}
where $t_m^{\texttt{f}}$ is the completion time of task $\omega_m$, and
$t^{\texttt{c}}$ is the scheduled time of the next communication event.
The objective in \eqref{ch05:eq:adaptive_objective} favors task sequences that
yield a high completion rate before the next communication event. In this way,
the formulation performs an implicit spatio-temporal filtering of candidate
task sequences without requiring an explicit model of the future task
distribution. The resulting metric balances task throughput against the time
needed to complete the selected tasks and restore team-wise communication.

\subsubsection{Overall algorithm}
\label{ch05:subsubsec:algorithm-description}

Let $\mathcal{T}$ denote the branch-and-bound search tree used by
Algorithm~\ref{ch05:alg:cocoplan}. Each vertex of $\mathcal{T}$ corresponds to
a partial assignment of detected tasks before the next communication event. A
node $\nu \in \mathcal{T}$ is associated with a partial collective plan defined as:
\[
J_\nu \triangleq
\bigl[
(\omega_{i_1}^1, \omega_{i_1}^2, \cdots, \mathbf{c}_{i_1}),
(\omega_{i_2}^1, \cdots, \mathbf{c}_{i_2}),
\cdots
\bigr],
\]
where each robot-specific task sequence is terminated by the next
communication event. Let $\Omega_\nu \subseteq \Omega_t$ denote the set of
tasks already inserted into $J_\nu$. For each node $\nu$, let $LB_\nu$ and
$UB_\nu$ denote, respectively, a lower bound and an upper bound on the
objective in \eqref{ch05:eq:adaptive_objective}.
Algorithm~\ref{ch05:alg:cocoplan} applies a branch-and-bound search to the
joint optimization of task assignment and communication scheduling for the full
robot team. For each node $\nu$, the objective is evaluated through
\eqref{ch05:eq:adaptive_objective}. The search maintains a lower bound
$LB_\nu$ and an upper bound $UB_\nu$ for every node. A node is pruned whenever
$UB_\nu \leq LB^\star$, where $LB^\star$ is the best lower bound found so far.
The procedure terminates when all nodes have been expanded or pruned, or when
the computation budget is exhausted. It then returns the best collective plan
$J^\star$ found within that budget.

\begin{algorithm}[!t]
    \caption{$\texttt{CoCoPlan}(\cdot)$}
    \label{ch05:alg:cocoplan}
    \SetAlgoLined
    \SetAlgoNlRelativeSize{-1}
    \DontPrintSemicolon
    \KwIn{Robot team $\mathcal{N}$, detected tasks $\Omega_t$, time budget $T'$}
    \KwOut{Best collective plan $J^\star$}

    Initialize $\mathcal{T} \gets \emptyset$ and max-heap $\mathcal{Q}$\;
    Create root node $\nu_0$ with $J_{\nu_0} \gets \emptyset$\;
    $(LB_{\nu_0}, J_{\nu_0})
    \gets \texttt{LowBound}(\nu_0, \Omega_t, \mathcal{N})$\;
    $UB_{\nu_0} \gets \texttt{UpBound}(\nu_0, \Omega_t, \mathcal{N})$\;
    $\mathcal{T} \gets \mathcal{T} \cup \{\nu_0\}$\;
    $\mathcal{Q}.\texttt{insert}(\nu_0, UB_{\nu_0})$\;
    $LB^\star \gets LB_{\nu_0}$ and $J^\star \gets J_{\nu_0}$\;
    $t_{\texttt{start}} \gets \texttt{current time}$\;

    \While{$\mathcal{Q} \neq \emptyset$ \textbf{and}
    $\texttt{current time} - t_{\texttt{start}} < T'$}{
        $\nu \gets \mathcal{Q}.\texttt{extractMax}()$\;

        \If{$UB_\nu > LB^\star$}{
            \If{$LB_\nu > LB^\star$}{
                $LB^\star \gets LB_\nu$ and $J^\star \gets J_\nu$\;
            }

            $\mathcal{F} \gets \texttt{GetTasks}(\nu, \Omega_t)$\;

            \ForEach{$\omega \in \mathcal{F}$}{
                $\nu_+ \gets \texttt{ExpNode}(\nu, \omega)$\;
                $(LB_{\nu_+}, J_{\nu_+})
                \gets \texttt{LowBound}(\nu_+, \Omega_t, \mathcal{N})$\;
                $UB_{\nu_+}
                \gets \texttt{UpBound}(\nu_+, \Omega_t, \mathcal{N})$\;
                $\mathcal{T} \gets \mathcal{T} \cup \{\nu_+\}$\;

                \If{$LB_{\nu_+} > LB^\star$}{
                    $LB^\star \gets LB_{\nu_+}$ and
                    $J^\star \gets J_{\nu_+}$\;
                }

                \If{$UB_{\nu_+} > LB^\star$}{
                    $\mathcal{Q}.\texttt{insert}(\nu_+, UB_{\nu_+})$\;
                }
            }
        }
    }
    \Return $J^\star$\;
\end{algorithm}
\setlength{\textfloatsep}{5pt}

(I) \emph{Initialization}.
Lines~1--8 initialize the search. The root node $\nu_0$ is created from the
empty collective plan. Its lower and upper bounds are computed by
$\texttt{LowBound}$ and $\texttt{UpBound}$, respectively. These values are
used to initialize the max-heap $\mathcal{Q}$ and the incumbent pair
$(LB^\star, J^\star)$. The timer $t_{\texttt{start}}$ is then recorded in
order to enforce the computation budget $T'$.

(II) \emph{Selection}.
At each iteration, the node with the largest upper bound is extracted from
$\mathcal{Q}$. This best-first rule prioritizes the nodes with the greatest
potential to improve the incumbent solution. If $UB_\nu \leq LB^\star$, then
$\nu$ and all of its descendants are discarded because they cannot improve the
current best plan.

(III) \emph{Expansion}.
For each surviving node, the routine $\texttt{GetTasks}$ identifies the set
$\mathcal{F}$ of feasible tasks that satisfy the temporal constraints
$\mathcal{C}_t$. The routine $\texttt{ExpNode}$ then generates child nodes by
assigning each task $\omega \in \mathcal{F}$ to eligible robots and inserting
it into their local plans before the next communication event. In this way,
task allocation and communication scheduling are coupled within a single search
structure.

(IV) \emph{Bounding}.
Each child node undergoes iterative bound refinement. At iteration $k$, a task
$\omega^{(k)} \in \Omega_t \setminus \Omega_{\nu^{(k)}}$ is selected and
assigned to an eligible robot group
$\mathcal{A} \in \mathcal{A}_k$ by solving:
\begin{equation}
\min_{\mathcal{A} \in \mathcal{A}_k}
\left\{
\max_{i \in \mathcal{A}}
T_{\text{travel}}(p_i^{\texttt{f}}, \omega^{(k)})
\right\},
\label{ch05:eq:task-assign}
\end{equation}
where $\mathcal{A}_k$ is the set of eligible robot groups at iteration $k$,
and $p_i^{\texttt{f}}$ is the current final planned position of robot $i$.
This step yields an updated plan $J^{(k)}$ for node $\nu^{(k)}$. Its
completion time is then evaluated in two phases. First, the non-communication
tasks are scheduled by minimizing the makespan
$\max_{\omega_m \in \Omega_{\nu^{(k)}}} t_m^{\texttt{f}}$ with a standard
linear-programming solver such as OR-Tools~\cite{glop}, while temporarily
ignoring communication decisions. Second, the routine $\texttt{ComOpt}$
optimizes the communication events and returns the corresponding communication
time $t_{(k)}^{\texttt{c}}$. The resulting candidate rate is
$b^{(k)} \triangleq |\Omega_{\nu^{(k)}}| / t_{(k)}^{\texttt{c}}$. Tasks are
added as long as $b^{(k)}$ improves, and the lower bound is then set to
$LB_{\nu_+} = \max_j b^{(j)}$. The upper bound $UB_{\nu_+}$ is computed by the
same procedure under the optimistic assumption of instantaneous robot motion.

 (V) \emph{Termination}.
The algorithm stops when the heap $\mathcal{Q}$ becomes empty or the time
budget $T'$ is exhausted. It then returns the incumbent plan $J^\star$ as the
best solution found. Since both the bounds and the incumbent are refined
progressively during the search, the method is anytime and can return a
feasible collective plan even under early termination.

The novelty of the proposed algorithm lies in its explicit co-optimization of
task assignment and communication scheduling within a single branch-and-bound
framework, rather than treating communication as a secondary feasibility check
after task allocation. By introducing the adaptive objective in
\eqref{ch05:eq:adaptive_objective}, the method evaluates collective plans
through the task-completion rate achieved before the next communication event,
which makes the coordination policy responsive to online task arrivals and
time-varying operating conditions. A further strength is that lower and upper
bounds are refined through both task scheduling and communication-event
optimization, so pruning is guided by a performance metric that reflects the
interaction between collaboration and information exchange. As a result, the
algorithm remains computationally effective through selective exploration of
promising plan branches, while preserving an anytime property that yields a
feasible collective plan even when the available computation time is limited.
























% ================================
% ================================
% ================================
\subsection{Optimization of fleet-wise communication events}
\label{ch05:subsec:comm-event}

Once a candidate non-communication task sequence has been fixed, the next fleet-wise communication event must be optimized explicitly. The two parts below define the event model and then describe the iterative refinement routine used to reduce waiting time while preserving connectivity.

\subsubsection{Fleet-wise communication events}
\label{ch05:subsubsec:comm-event-opt}

Given the non-communication plan $J_{\nu_+}$ obtained in the previous step,
the remaining task is to determine the next team-wise communication event.
This step is necessary because the adaptive coordination framework in
Section~\ref{ch05:subsec:coord-framework} evaluates a collective plan only up
to the next communication time. Consequently, after the task sequence before
that event has been determined, the corresponding communication positions and
the synchronized communication time must also be optimized. The purpose of this
subsection is therefore to refine the communication event associated with the
candidate collective plan, so that the fleet can re-establish global
connectivity with small waiting time.
Let $t_i^{\texttt{f}}$ denote the completion time of the last
non-communication task assigned to robot $i$ in the local plan $\tau_i$. Let
$p_i^{\texttt{f}} \in \mathcal{W}$ denote the corresponding robot position at
time $t_i^{\texttt{f}}$. The subsequent communication event must occur at a
time $t^{\texttt{c}} > \max_{i \in \mathcal{N}} t_i^{\texttt{f}}$, with robots
located at positions $\{p_i^{\texttt{c}}\}_{i \in \mathcal{N}}$, where
$p_i^{\texttt{c}} \in \mathcal{W}_i$, such that the induced communication graph
satisfies the connectivity requirement in \eqref{ch05:eq:comm-graph}. This
means that the communication event must be feasible in both motion and
communication terms. Each robot must be able to reach its assigned
communication position after completing its last task, and the resulting
fleet-wide configuration must support connected communication at the common
event time.

The communication-event optimization is formulated so as to reduce the largest
idle delay between the completion of the last assigned task and the next
communication event. In this way, the fleet avoids configurations in which one
or more robots finish early and then wait excessively long for the remaining
robots. The resulting optimization problem is given by:
\begin{equation}
\begin{aligned}
\min_{t^{\texttt{c}},\, \{p_i^{\texttt{c}}\}} \quad &
\max_{i \in \mathcal{N}} \bigl(t^{\texttt{c}} - t_i^{\texttt{f}}\bigr), \\
\text{s.t.} \quad &
t^{\texttt{c}} \geq t_i^{\texttt{f}} +
T_{\text{travel}}(p_i^{\texttt{f}}, p_i^{\texttt{c}}),
\quad \forall i \in \mathcal{N}, \\
&
G(t^{\texttt{c}}) \in \mathcal{G}_{\texttt{c}},
\end{aligned}
\label{ch05:eq:comm-opt}
\end{equation}
where the first constraint in \eqref{ch05:eq:comm-opt} guarantees that robot
$i$ has sufficient time to move from its last task position
$p_i^{\texttt{f}}$ to its assigned communication position
$p_i^{\texttt{c}}$, and the second enforces that the fleet is connected at the
communication time. Hence, \eqref{ch05:eq:comm-opt} couples two requirements
that are both indispensable to the subsequent replanning stage: temporal
synchronization and fleet-wise connectivity.
The objective in \eqref{ch05:eq:comm-opt} is consistent with the rate-based
criterion introduced in \eqref{ch05:eq:adaptive_objective}. Indeed, once the
non-communication portion of the plan has been fixed, reducing the communication
time directly improves the denominator of the adaptive objective. For this
reason, communication-event optimization should not be viewed as a secondary
post-processing step. It is part of the same coordination objective, because
the time at which communication is restored directly influences the effective
task-completion rate of the collective plan.

\subsubsection{Optimization by iterative refinement}
\label{ch05:subsubsec:comm-event-algorithm}

\begin{algorithm}[!t]
    \caption{$\texttt{ComOpt}(\cdot)$}
    \label{ch05:alg:comopt}
    \SetAlgoLined
    \DontPrintSemicolon
    \KwIn{Robot team $\mathcal{N}$, last-task data
    $\{(t_i^{\texttt{f}}, p_i^{\texttt{f}})\}$, time budget $T'$,
    thresholds $\delta$ and $\Delta$}
    \KwOut{Communication event $\mathbf{C}_t$}

    $i_\ell \gets \arg\max_{i \in \mathcal{N}} t_i^{\texttt{f}}$\;
    $p_0 \gets p_{i_\ell}^{\texttt{f}}$\;
    $t^{\texttt{c}} \gets t_{i_\ell}^{\texttt{f}}$,
    $t^{\texttt{b}} \gets t^{\texttt{c}}$\;
    $\mathbf{C}_t \gets \{(t^{\texttt{b}}, p_0) \mid i \in \mathcal{N}\}$\;
    $t_{\texttt{start}} \gets \texttt{current time}$\;

    \While{$\texttt{current time} - t_{\texttt{start}} < T'$
    \textbf{and} $|t^{\texttt{c}} - t^{\texttt{b}}| \geq \Delta$}{
        $i_{\texttt{e}} \gets
        \arg\min_{i \in \mathcal{N}}
        \bigl(t_i^{\texttt{f}} +
        \texttt{A}^\star(p_i^{\texttt{f}}, p_0)\bigr)$\;

        $j_n \gets
        \arg\min_{j \neq i_{\texttt{e}}}
        \texttt{A}^\star(p_{i_{\texttt{e}}}^{\texttt{f}}, p_j^{\texttt{c}})$\;

        $p^{\texttt{c}} \gets
        \texttt{SelCom}(p_{i_{\texttt{e}}}^{\texttt{f}}, p_{j_n}^{\texttt{c}})$\;

        $t_{\texttt{a}} \gets
        t_{i_{\texttt{e}}}^{\texttt{f}} +
        \texttt{A}^\star(p_{i_{\texttt{e}}}^{\texttt{f}}, p^{\texttt{c}})$\;

        $t_+^{\texttt{c}} \gets
        \max \Bigl(
        t_{\texttt{a}},
        \max_{k \neq i_{\texttt{e}}}
        \bigl(t_k^{\texttt{f}} +
        \texttt{A}^\star(p_k^{\texttt{f}}, p_k^{\texttt{c}})\bigr)
        \Bigr)$\;

        \If{$Q_{i_{\texttt{e}}j_n}(t_+^{\texttt{c}}) > \delta$
        \textbf{and} $t_+^{\texttt{c}} < t^{\texttt{b}}$}{
            $p_{i_{\texttt{e}}}^{\texttt{c}} \gets p^{\texttt{c}}$\;
            \If{$|t^{\texttt{b}} - t_+^{\texttt{c}}| < \Delta$}{
                $t^{\texttt{c}} \gets t_+^{\texttt{c}}$,
                $t^{\texttt{b}} \gets t_+^{\texttt{c}}$\;
                \textbf{break}\;
            }
            $t^{\texttt{c}} \gets t_+^{\texttt{c}}$,
            $t^{\texttt{b}} \gets t_+^{\texttt{c}}$\;
        }
    }
    \Return $\mathbf{C}_t =
    \{(t^{\texttt{b}}, p_i^{\texttt{c}}) \mid i \in \mathcal{N}\}$\;
\end{algorithm}

Algorithm~\ref{ch05:alg:comopt} addresses the optimization problem in
\eqref{ch05:eq:comm-opt} through iterative refinement of the communication
configuration. Rather than enumerating all feasible communication positions,
the procedure starts from a simple reference configuration and improves it
gradually. This design is appropriate in the present framework because the
communication-event optimization is invoked repeatedly inside the higher-level
coordination loop. A lightweight refinement procedure is therefore preferable
to an exhaustive global search.

(I) \emph{Initialization}.
The procedure begins with the robot that completes its last non-communication
task the latest. Let this robot be $i_\ell$. Since no communication event can
occur before robot $i_\ell$ finishes its assigned task, this robot determines
the earliest feasible reference time for the next fleet-wise communication
event. Its final task position $p_{i_\ell}^{\texttt{f}}$ is used as the
initial communication location $p_0$, and the synchronized communication time
is initialized by
$t^{\texttt{c}} = t^{\texttt{b}} = t_{i_\ell}^{\texttt{f}}$. This yields an
initial event in which all robots are directed toward the same communication
location. Although this initialization is conservative, it provides a feasible
baseline from which subsequent refinements can reduce the overall waiting time.

(II) \emph{Selection of the robot to be refined}.
Within the computation budget $T'$, the algorithm improves the initial
configuration iteratively until either convergence is reached or the available
time is exhausted. At each iteration, the earliest-arriving robot
$i_{\texttt{e}}$ is selected by minimizing
$t_i^{\texttt{f}} + \texttt{A}^\star(p_i^{\texttt{f}}, p_0)$. This robot can
reach the current communication location sooner than the others and is thus
the most natural candidate for relocation. Intuitively, robots that would
otherwise arrive much earlier than the synchronized communication time are the
ones contributing most strongly to avoidable idle waiting.

(III) \emph{Construction of a candidate communication position}.
After robot $i_{\texttt{e}}$ is selected, the algorithm identifies a
neighboring anchor robot $j_n$ according to the shortest
$\texttt{A}^\star$ path from $p_{i_{\texttt{e}}}^{\texttt{f}}$ to the current
communication position $p_{j_n}^{\texttt{c}}$. This step refines the
configuration locally while preserving the possibility of connected
communication. A candidate communication position $p^{\texttt{c}}$ is then
generated between $p_{i_{\texttt{e}}}^{\texttt{f}}$ and $p_{j_n}^{\texttt{c}}$
through the routine $\texttt{SelCom}$. This candidate moves the communication
position of robot $i_{\texttt{e}}$ toward a more balanced configuration in
time, while remaining compatible with the underlying communication geometry.

(IV) \emph{Evaluation and acceptance of the candidate}.
Once the candidate position has been obtained, the algorithm computes the
arrival time $t_{\texttt{a}}$ of robot $i_{\texttt{e}}$ at the new position.
The synchronized communication time is then updated to $t_+^{\texttt{c}}$,
which is defined as the largest arrival time among all robots under the
current candidate configuration. Hence, $t_+^{\texttt{c}}$ is the earliest
fleet-wise communication time compatible with the updated positions. The
candidate is accepted only if two conditions hold. First, the communication
quality between the updated position of robot $i_{\texttt{e}}$ and the anchor
position of robot $j_n$ exceeds the threshold $\delta$. Second, the new
synchronized time $t_+^{\texttt{c}}$ is strictly smaller than the current best
time $t^{\texttt{b}}$. Therefore, an update is retained only when it improves
temporal efficiency without violating communication feasibility.

(V) \emph{Termination and output}.
The refinement process terminates when the improvement falls below the
tolerance $\Delta$ or when the computation budget is exhausted. The output of
Algorithm~\ref{ch05:alg:comopt} is the communication event $\mathbf{C}_t$,
which consists of the communication positions
$\{p_i^{\texttt{c}}\}_{i \in \mathcal{N}}$ and the synchronized communication
time $t^{\texttt{b}}$. By construction, this event satisfies the connectivity
requirement in \eqref{ch05:eq:comm-graph} and reduces the maximum delay
$\max_{i \in \mathcal{N}} (t^{\texttt{b}} - t_i^{\texttt{f}})$ in accordance
with \eqref{ch05:eq:comm-opt}. The strength of the procedure lies in the fact
that it replaces a global search over communication configurations by a
sequence of inexpensive local refinements guided jointly by task-completion
times, motion costs, and communication feasibility.













% ================================
\subsection{Overall analysis}
\label{ch05:subsec:online}

The preceding constructions define a closed online coordination loop rather than a one-shot offline optimizer. The two parts below discuss how the scheme is executed during deployment and how its computational burden scales with team size and detected tasks.

\subsubsection{Online execution}
\label{ch05:subsubsec:online-execution}

To accommodate dynamically emerging tasks, the coordination framework is
executed online. Robots discover new tasks while moving in the workspace, and the detected task set is
updated incrementally as $\Omega_t$. During execution, each robot follows its
current local plan $\tau_i$ and proceeds toward either the next assigned task
or the next communication location. In this way, planning and execution remain
interleaved throughout the mission, rather than being separated into a single
offline synthesis phase and a purely reactive deployment phase.
To absorb execution uncertainty, simple synchronization rules are imposed at
the task and communication levels. For each assigned collaborative task, the
corresponding robots wait at the task region until all required robots have
arrived, after which task execution begins. Likewise, at each communication
event, all robots wait until every robot has reached its designated
communication position. As a result, the connectivity condition
\eqref{ch05:eq:comm-graph} is enforced at the communication time even when
travel times or execution durations deviate from their nominal values. These
waiting rules introduce a basic robustness mechanism without modifying the
underlying coordination structure.
At every communication event, the robots exchange information about newly
detected tasks, current execution progress, and future plan updates. A
designated robot then aggregates the locally discovered tasks into the global
set $\Omega_t$ and invokes Algorithm~\ref{ch05:alg:cocoplan} to compute an
updated collective plan together with the next fleet-wise communication event.
The robots subsequently execute the revised plan until the next communication
event is reached, where the same update procedure is repeated. The overall
execution logic is therefore cyclic: task execution, task discovery,
communication, and replanning are repeated in succession as new information
becomes available.

\subsubsection{Complexity analysis}
\label{ch05:subsubsec:complexity}

The computational complexity of
Algorithm~\ref{ch05:alg:cocoplan} is dominated by the branch-and-bound search.
Suppose that $m$ tasks have been detected and that the fleet contains $N$
robots. In the worst case, a single expansion step may generate up to
$\mathcal{O}(m \cdot N!)$ child nodes, since each detected task may be assigned
to multiple robot orderings or coalition patterns. Repeating this process over
a search depth of at most $m$ yields a worst-case tree size of
$\mathcal{O}((m \cdot N!)^m)$. For each node, the dominant cost remains the
evaluation of candidate task expansions and the associated bound computations,
which again scale as $\mathcal{O}(m \cdot N!)$. Consequently, the overall
worst-case complexity remains exponential in the number of detected tasks.
This worst-case bound is conservative and mainly reflects the combinatorial
nature of online multi-robot task assignment under temporal and communication
constraints. In practice, the effective search space is typically much smaller.
The upper and lower bounds computed at each node allow large portions of the
tree to be discarded early, especially when the adaptive objective in
\eqref{ch05:eq:adaptive_objective} quickly separates promising task sequences
from unproductive ones. The branch-and-bound structure is therefore important
not only for correctness, but also for computational tractability within a
finite replanning budget.

\begin{lemma}
\label{ch05:lem:alg2_graph_connect}
Algorithm~\ref{ch05:alg:comopt} returns communication locations such that the
communication graph $G(t^{\texttt{c}})$ is connected.
\end{lemma}

\begin{proof}
The initial communication event places all robots at the same communication
location $p_0$. Hence, the induced communication graph is connected at the
initialization step. During each refinement step, the algorithm updates the
communication location of only one robot. The new location is accepted only if
it preserves a valid communication link to an already selected anchor location.
Therefore, the updated robot remains connected to the set of robots whose
communication locations have already been fixed. Repeating this argument over
all accepted refinements shows inductively that connectivity is preserved
throughout the algorithm. Consequently, the final communication graph at
$t^{\texttt{c}}$ is connected.
\end{proof}

\begin{theorem}
\label{ch05:thm:cocoplan_opt}
Assume that communication events recur indefinitely and that, at every
communication time, Algorithm~\ref{ch05:alg:cocoplan} is solved to optimality
for the adaptive objective in \eqref{ch05:eq:adaptive_objective}. Then, on each
communication cycle $[t_k^{\texttt{c}}, t_{k+1}^{\texttt{c}})$, the returned
plan is feasible with respect to \eqref{ch05:eq:task-constraints} and
\eqref{ch05:eq:comm-graph}, and locally maximizes the task-completion rate over
that cycle. Moreover, \eqref{ch05:eq:adaptive_objective} provides a local
approximation of the global infinite-horizon objective
\eqref{ch05:eq:objective}. Under stationary task-arrival and ergodicity
assumptions, the two optima coincide.
\end{theorem}

\begin{proof}
Feasibility on each communication cycle follows from the construction of the
algorithm. Algorithm~\ref{ch05:alg:cocoplan} enforces the temporal constraints
in \eqref{ch05:eq:task-constraints} during task expansion and bound
computation. It also invokes Algorithm~\ref{ch05:alg:comopt} to determine the
next communication event, and Lemma~\ref{ch05:lem:alg2_graph_connect} shows
that this event satisfies the connectivity requirement
\eqref{ch05:eq:comm-graph}. Therefore, every returned cycle plan is feasible
with respect to both task ordering and team-wise communication.
Local optimality follows from the assumption that
Algorithm~\ref{ch05:alg:cocoplan} is solved to optimality for
\eqref{ch05:eq:adaptive_objective} at each communication event. Under that
assumption, the returned plan maximizes the number of tasks completed per unit
time before the next communication event, and therefore yields the largest
adaptive completion rate over the current cycle.
The relation to the infinite-horizon objective
\eqref{ch05:eq:objective} is then immediate at the approximation level. The
adaptive objective evaluates the same performance quantity as
\eqref{ch05:eq:objective}, but only over a single communication cycle rather
than over the full mission horizon. If task arrivals are stationary and the
resulting communication cycles are ergodic, then the long-run average task
completion rate equals the time average of the per-cycle completion rates. In
that case, optimizing each cycle separately is equivalent to optimizing the
infinite-horizon average objective, and the local optimum induced by
\eqref{ch05:eq:adaptive_objective} coincides with the global optimum in
\eqref{ch05:eq:objective}.
\end{proof}


This section showed that communication should be treated as a design
variable of the coordination process rather than as a passive execution
constraint. Collaborative-task assignment, fleet-wise communication
events, and online replanning were integrated within a single
optimization loop, yielding cyclewise feasibility and a principled
approximation of the long-run objective. The next section studies a
complementary setting in which intermittent connectivity is not merely
scheduled, but structurally embedded into the exploration process.
